\documentclass[../../main.tex]{subfiles}
\begin{document}
\begin{oframed}
  \subfile{problem}
\end{oframed}

{\bf [解答]}

与えられた関数
\begin{align}
  f(\theta) = \cos 4\theta - 4\sin^2\theta
\end{align}
を$\theta$で微分して
\begin{align}
  f'(\theta)
  &=-4\sin4\theta-8\sin\theta\cos\theta \\
  &=-8\sin2\theta\cos2\theta-8\sin\theta\cos\theta\\
  &=-8\cos\theta\sin\theta(2\cos2\theta+1)\\
  &=-8\cos\theta\sin\theta(4\cos^2\theta-1)
\end{align}
だから，増減表は以下のようになる．

\begin{table}[htb]
  \centering
  \caption{$0\leqq\theta\leqq \dfrac{3\pi}{4}$における$f(\theta)$の増減表}
  \begin{tabular}{c|ccccccccc}
    $\theta$ & $0$ & & $\pi/3$ & & $\pi/2$ & & $2\pi/3$ & & $3\pi/4$ \\
    \hline
    $f'$ & & $-$ & $0$ & $+$ & $0$ & $-$ & $0$ & $+$ & \\
    \hline
    $f$ & $1$ & $\searrow$ & $-\dfrac{7}{2}$ & $\nearrow$ & $-3$ & $\searrow$ & $-\dfrac{7}{2}$ & $\nearrow$ & $-3$ \\
  \end{tabular}
\end{table}

したがって求める最大・最小値は
\begin{align}
  \begin{cases}
    \max f(\theta) = f(0) = 1 \\
    \min f(\theta) = f(\pi/3) = f(2\pi/3) = -\dfrac{7}{2}
  \end{cases}
\end{align}
である．

\newpage

\end{document}
